% !TeX root = ncaa-manuscript-en.tex
% English Introduction revised from the author's Korean draft on 2026-10-17.
% The preceding version is preserved in 01-introduction-v02-en.tex.
\section{Introduction}\label{sec:introduction}

Vision-language models (VLMs) and vision-language-action (VLA) models connect
visual observations and language information to action selection or trajectory
prediction, with autonomous driving as a prominent application. Recent work uses supervised fine-tuning and
reinforcement learning to improve language reasoning and action prediction
together~\cite{r1}. Other proposals integrate future-scene modeling into planning
or maintain language-based memories of past observations to support decisions
over time~\cite{Wang2026DriveWLA,Yan2026ADMemo}. Evaluation also extends from
individual prediction performance to closed-loop driving and behavioral changes
under external instructions~\cite{Jia2026Bench2DriveVL,You2026Language2Drive}.

Despite these advances, reliably connecting a model's judgment to the required
behavior remains difficult. Drive the Thoughts reports only 74\% agreement
between predicted trajectories and reasoning traces judged reliable~\cite{r18}.
Language2Drive reports that models can succeed under one instruction yet fail
under a feasible alternative when the initial scene is held
fixed~\cite{You2026Language2Drive}. Generating an explanation or driving
successfully under one condition therefore does not establish the ability to
adjust behavior when external conditions change. We focus on using externally
supplied behavioral rules as conditions for action selection. A supplied
behavioral rule states a condition to apply in a particular case; its source
may be a traffic regulation, an operator instruction, or an experimental
policy. The same scene and action candidates can admit different actions
under different rules.

Using such rules is distinct from explaining the reason for a driving
decision. Driving models that reason in language express how an observed
situation leads to a decision. Alpamayo-R1's Chain of Causation (CoC) is one
example: it learns causal relationships between observations and decisions
together with trajectory prediction~\cite{r1}. Such reasoning need not fully
specify an external rule's applicability conditions and numerical parameters
as a case-specific behavioral constraint. Here, a behavioral constraint makes
a supplied rule concrete through applicability conditions, action restrictions,
timing, and release conditions, so that action admissibility can be assessed.
The scene and causal explanation can remain valid while different speed
bounds, waiting durations, or stop-release conditions change the admissible
actions. An action consistent with its explanation can therefore violate a
separately supplied constraint.

Behavioral constraints determine which actions are admissible, but do not by
themselves rank the admissible actions. For example, always choosing the longer
wait under two different waiting rules may satisfy both rules while missing an
opportunity for earlier permitted progress. Constraint compliance and
preferences among admissible actions must therefore be distinguished. In our
candidate-selection task, the goal is to satisfy the applicable constraints
and select the admissible candidate with the highest supplied progress score.
This score is a synthetic value assigned by the task generator; it is not a
direct measurement of distance traveled or long-term driving efficiency.

Learning this selection requires training examples to make the relationship
among scenes, applicable rules, and actions explicit. It also requires checking
whether constraints assembled from observations, supplied rules, and modeling
assumptions express the intended behavioral conditions, and assessing whether
the model uses them in its choices. Including a rule in the input alone does
not establish such use. Our preliminary Alpamayo-R1 diagnostics, without
additional training, produced predictions violating explicitly supplied speed
and position-hold rules on sampled PhysicalAI~\cite{nvidia2025physicalai} and
converted nuScenes~\cite{caesar2020nuscenes} inputs. In a separate controlled
candidate-selection study, we therefore test whether including constraints in
training improves constraint compliance and objective-based candidate selection
when every model receives the same rule at evaluation.

Behavioral constraints supplied for learning require rigorous validation to
determine whether their combination satisfies the specified safety requirements
under the stated assumptions. Combining behavioral conditions can restrict a
necessary safety response. The 2018 fatal Uber collision in Tempe illustrates
this concern. The National Transportation Safety Board (NTSB) reported emergency
recognition 1.2~s before impact, followed by 1~s of braking suppression intended
to reduce false-alarm responses. Only about 0.2~s remained afterward. The system
also lacked maximum allowed braking solely for mitigation when collision
avoidance exceeded its braking limits. The NTSB identified increased safety
risk from this design~\cite{NTSB2019Tempe}. This illustrates why suppression
timing must be checked against collision-response requirements.

In the 2023 Cruise pedestrian collision, incorrect lateral-collision
classification led to a pullover after an initial stop, dragging the
pedestrian~\cite{Cruise2023Recall23E086}. This case highlights the need to
examine the conditions governing transitions between remaining stationary and
moving. Rule 129 of the United Kingdom's Highway Code generally prohibits
crossing double white centre lines when the
line on the driver's side is solid, but allows crossing when necessary to pass
a stationary vehicle, provided the road is clear~\cite{DfT2015HighwayCode129}.
This rule illustrates why specifications must preserve exceptions in the
original rules: a constraint that omits this exception could block a lawful
manoeuvre and prolong unnecessary waiting. Our learning experiments use
single constraints, while a separate logical diagnostic combines conditions
drawn from these and related documents, including the Tempe braking
suppression and a restricted encoding of the Rule 129 exception, under
explicitly authored states. It examines the joint feasibility of the stated
conditions, rather than reconstructing the accidents or interpreting the
regulations in full.

These examples motivate treating behavioral constraints as checkable objects
with explicit evidence and applicability. Supplying them as natural-language
learning inputs raises three problems. First, the evidence and roles of
observed facts, supplied rules, and modeling assumptions must be explicit, so
that observation errors, rule misapplication, and unsuitable assumptions can
be examined separately. Second, the specification must express the intended
action restrictions and release conditions. For the learning specifications
evaluated here, satisfiability and candidate admissibility must be checked
under the supplied inputs and assumptions. When several supplied conditions apply, joint satisfiability must also be
assessed under an explicit shared context. Third, mechanical checking
requires a formal representation, whereas the model learns from natural
language. Meaning must be preserved when the specification is rendered as a
sentence, and model behavior must be evaluated separately: correspondence
between a sentence and its specification does not establish constraint compliance.

To address these problems, we propose GuardSynth. It first binds each
constraint field to an observation, a supplied rule, or a modeling assumption,
separating their roles. It represents applicability, action restrictions,
persistence and release, timing, and evidence in Evidence-Bound Lifecycle
Contracts (EBLC). Supported specifications are lowered to Satisfiability Modulo
Theories (SMT) queries for satisfiability and candidate-admissibility checks
within a restricted logical model. A separate composition diagnostic checks
source conditions jointly under supplied contexts. GuardSynth then generates Controlled Natural Language
(CNL) from the specification and appends it to the original CoC. For supported
fixed grammars, parsing the generated sentence recovers fields for comparison
with the specification. Our evaluation concerns the supported specifications
and synthetic candidate-selection tasks used in this process.

The evaluation separates the behavioral effect of training with constraints
from specification checkability. Constraint violation rate is the proportion
of all predictions that select a prohibited candidate. The
\emph{preferred-selection score} is the proportion that select the most
preferred admissible candidate under the specified objective. In our tasks,
the preferred candidate has the highest supplied progress score among
admissible candidates; selecting its label receives a score of one, and every
other output receives zero. The mean is reported as a percentage. Rule-pair
accuracy is the proportion of pairs for which both reference labels are
selected when the scene and candidates are fixed and only the rule changes.
These metrics distinguish constraint satisfaction from objective-based
selection within the supplied candidate set.

We evaluate fresh Qwen3-VL-2B-Instruct adapters~\cite{qwen2025qwen3vl2b} with
matched initialization and the same CNL at evaluation over three training
seeds. Training with CNL under cross-entropy loss alone yields three-seed mean
constraint violation rates of at most 1.56\% and rule-pair accuracy of at least
86.98\% across pedestrian entry and the two position-hold representations.
This controlled comparison supports including case-specific constraints in
training data as well as supplying them at inference for the evaluated
adapters. Across six behavioral tasks, both natural-language guards (P1) and
EBLC-derived CNL (P2) improve selection of the preferred candidate under the
supplied rule; this comparison changes rule availability as well as training
inputs.

A separate fault-injection evaluation of the two specifications used to
generate learning inputs detects all 288 errors of six supported types and
accepts all 48 valid cases. All 48 jointly wrong evidence--specification
controls remain undetected: the evidence record and specification are changed
to the same wrong value. This identifies a boundary between internal
correspondence checks and independent confirmation of the evidence itself.
Evidence correspondence, sentence correspondence, and boundary probes for
logic--rule comparison were added after training for this evaluation; their
results are interpreted separately from the causes of behavioral improvement.

A separate composition diagnostic combines manually formalized conditions
from official documents with explicit modeling assumptions and authored
states, candidates, or test objectives. Among 176 configurations, it identifies
58 incompatible combinations, including 12 three-way conflicts for which
every pair is satisfiable. All 118 compatible combinations are retained. These results concern joint logical
feasibility under explicit contexts, rather than VLM behavior or a general
vehicle-safety guarantee.

Our contributions are:
\begin{enumerate}
\item A procedure for constructing case-specific behavioral constraints tied
to their evidence. It separates observations, supplied rules, and modeling
assumptions and links them to EBLC specifications with applicability, action
restrictions, persistence, and release conditions, making the components and
evidence of learning constraints explicit objects of checking.
\item The GuardSynth development path connecting formal specification checks
to natural-language learning inputs. It lowers supported specifications to SMT
for satisfiability and candidate-admissibility checks within a restricted model,
appends specification-derived CNL to the original CoC, and checks
specification--sentence correspondence by parsing supported grammars. The same
Core/SMT compiler checks bounded source combinations under shared contexts in a
separate diagnostic that does not produce learning inputs.
\item An experimental evaluation separating the effect of training with
constraints from specification checkability. Matched initialization and
evaluation inputs test improvements in constraint compliance and
objective-based candidate selection, while separate fault injection identifies
complementary checking layers and their detection limits. A separate
composition diagnostic compares individual, pairwise, and joint consistency
checks for supplied sources. Model behavior and
the correctness of learning-constraint construction are assessed separately.
\end{enumerate}

Section~\ref{sec:related-work} reviews related work, and
Section~\ref{sec:method} presents the problem setting and GuardSynth.
Sections~\ref{sec:experimental-setup} and~\ref{sec:results} describe the
experimental setup and results. Section~\ref{sec:discussion} discusses
implications and limitations, and Section~\ref{sec:conclusion} concludes.
